\documentclass[11pt,a4paper]{article}

% ------------------------------------------------------------------
% Packages (all available on Overleaf / ShareLaTeX)
% ------------------------------------------------------------------
\usepackage[utf8]{inputenc}
\usepackage[T1]{fontenc}
\usepackage{lmodern}
\usepackage[margin=2.3cm]{geometry}
\usepackage{amsmath,amssymb,amsthm,bm,mathtools}
\usepackage{booktabs,array,longtable}
\usepackage{enumitem}
\usepackage[most]{tcolorbox}
\usepackage{algorithm}
\usepackage{algpseudocode}
\usepackage[colorlinks=true,linkcolor=blue!50!black,citecolor=blue!50!black]{hyperref}

% ------------------------------------------------------------------
% Notation
% ------------------------------------------------------------------
\newcommand{\C}{\mathbb{C}}                     % homogenized stiffness (known)
\newcommand{\Cb}{\bar{\mathbb{C}}}              % relative-strain stiffness
\newcommand{\Ch}{\hat{\mathbb{C}}}              % coupling stiffness
\newcommand{\GG}{\mathbb{G}}                    % curvature matrix
\newcommand{\Dm}{\mathbf{D}}                    % phase stiffness
\newcommand{\B}{\mathbb{B}}                     % local energy block
\newcommand{\eps}{\bm{\varepsilon}}
\newcommand{\epsM}{\bm{\varepsilon}_{M}}
\newcommand{\bpsi}{\bm{\psi}}
\newcommand{\bgam}{\bm{\gamma}}
\newcommand{\bkap}{\bm{\kappa}}
\newcommand{\bsig}{\bm{\sigma}}
\newcommand{\bmu}{\bm{\mu}}
\newcommand{\bx}{\bm{x}}
\newcommand{\bX}{\bm{X}}
\newcommand{\bu}{\bm{u}}
\newcommand{\bH}{\bm{H}}
\newcommand{\bG}{\bm{G}}
\newcommand{\bK}{\bm{K}}
\newcommand{\bP}{\bm{P}}
\newcommand{\bb}{\bm{b}}
\newcommand{\bth}{\bm{\theta}}
\newcommand{\sym}{\operatorname{sym}}
\newcommand{\dd}{\,\mathrm{d}}
\newcommand{\ve}[1]{\{#1\}_{\varepsilon}}      % strain-like Voigt vector
\newcommand{\vs}[1]{\{#1\}_{\sigma}}           % stress-like Voigt vector
\newcommand{\Zset}{\mathcal{Z}}
\newcommand{\Nset}{\mathcal{N}}
\newcommand{\Kset}{\mathcal{K}}
\newcommand{\code}[1]{\texttt{#1}}

\theoremstyle{plain}
\newtheorem{proposition}{Proposition}
\newtheorem{lemma}{Lemma}
\theoremstyle{remark}
\newtheorem{remark}{Remark}

% ------------------------------------------------------------------
% "Delivers" box for every step
% ------------------------------------------------------------------
\newtcolorbox{delivers}{colback=green!4!white, colframe=green!40!black,
  boxrule=0.6pt, arc=2pt, left=4pt, right=4pt, top=2pt, bottom=2pt,
  fonttitle=\bfseries\small, title=Delivers, fontupper=\small}
\newtcolorbox{codebox}{colback=blue!3!white, colframe=blue!40!black,
  boxrule=0.4pt, arc=2pt, left=4pt, right=4pt, top=1pt, bottom=1pt,
  fontupper=\small}
\newcommand{\implementedin}[1]{\begin{codebox}\textbf{Implemented in:} #1\end{codebox}}

\title{Identification of the Micromorphic Stiffness Matrices\\
$\Cb$, $\Ch$ and the Curvature Modulus $G$ from Classical FE Data\\[4pt]
\large Mathematical basis of \code{main\_identify\_CbarChatG.m}, \code{dataGeneration.m},\\
\code{homogenize2D\_PBC.m}, \code{smoothData.m}, \code{identifyCbarChatG.m}, \code{plotAterms.m}}
\author{}
\date{}

\begin{document}
\maketitle

\begin{abstract}
The document states, stage by stage and step by step, every mathematical operation performed by
the identification code. Classical plane-strain finite-element (FE) solutions of an $N\times N$
matrix--inclusion arrangement are computed under kinematic uniform boundary conditions with
prescribed first and second displacement gradients, averaged over a full period of inclusion
translations, and reduced to domain integrals of the micro-deformation $\bpsi$, the relative strain
$\bgam$ and the micro-curvature $\bkap$, together with the heterogeneous strain energy, the
integrated stress and the integrated first stress moment. After shape-constrained smoothing over
$N$, the matrices $\Cb$, $\Ch$ and the curvature modulus $G$ of a micromorphic strain energy are
identified by a convex, positive-definiteness-constrained least-squares problem in which every
parameter is a positive, non-increasing function of $N$.
\end{abstract}

\tableofcontents
\bigskip

\begin{table}[h]
\centering\small
\begin{tabular}{@{}lll@{}}
\toprule
Stage & Content & Code\\
\midrule
0 & Setting, conventions, micromorphic model & --\\
1 & Load cases and kinematic boundary conditions & \code{dataGeneration}\\
2 & Homogenized stiffness $\C$ (periodic unit cell) & \code{homogenize2D\_PBC}\\
3 & Finite-element discretization & \code{dataGeneration}\\
4 & Evaluation grid, quadrature, regions & \code{dataGeneration}\\
5 & Macroscopic reference field $\epsM$ & \code{dataGeneration}\\
6 & Translation averaging and FE stress integrals & \code{dataGeneration}\\
7 & Relative strain and micro-curvature & \code{dataGeneration}\\
8 & Integral data $A_1\dots A_{19}$, $I_1\dots I_{24}$ & \code{dataGeneration}\\
9 & Smoothing over $N$ & \code{smoothData}\\
10 & Identification equations & \code{identifyCbarChatG}\\
11 & Solution: monotone, energy-positive-definite (default) & \code{identifyCbarChatG}\\
12 & Solution: free per $N$ / global (comparison) & \code{identifyCbarChatG}\\
13 & Outputs and diagnostic plots & \code{identifyCbarChatG}, \code{plotAterms}\\
14 & Consistency properties and checks & all\\
\bottomrule
\end{tabular}
\caption{Stages of the pipeline, in the order in which \code{main\_identify\_CbarChatG.m} executes them.}
\end{table}

% ==================================================================
\section{Stage 0 -- Setting, conventions and the micromorphic model}
\label{sec:stage0}
% ==================================================================

\subsection*{Step 0.1 -- Specimen, unit cells and phases}
The specimen is $\Omega=[0,L_x]\times[0,L_y]$ ($L_x=L_y=L=1$ by default) of thickness $t$ and
volume $V=L_xL_yt$, divided into $N\times N$ square unit cells of size
\begin{equation}
  \ell=\frac{L}{N},\qquad \ell^2(N)=\Big(\frac{L_x}{N}\Big)^2 ,
  \label{eq:ell}
\end{equation}
each containing one circular inclusion of radius $R=r\,\ell$ ($r=6/19$). Coordinates relative to
the centre $\bX_c=(L_x/2,L_y/2)$ are $\bx=\bX-\bX_c$. Both phases are isotropic, linear elastic,
in plane strain:
\begin{equation}
  \Dm(E,\nu)=\frac{E}{(1+\nu)(1-2\nu)}
  \begin{bmatrix}1-\nu&\nu&0\\ \nu&1-\nu&0\\ 0&0&\tfrac{1-2\nu}{2}\end{bmatrix},
  \qquad
  \Dm(\bX)=\begin{cases}\Dm_m=\Dm(E_m,\nu_m),&\bX\in\Omega_m,\\ \Dm_i=\Dm(E_i,\nu_i),&\bX\in\Omega_i.\end{cases}
  \label{eq:Dphase}
\end{equation}
\begin{delivers}
Geometry $(L_x,L_y,t,V,N,\ell,R)$ and phase stiffnesses $\Dm_m,\Dm_i$.
\end{delivers}

\subsection*{Step 0.2 -- Tensor and Voigt conventions}
Strain-type fields $\eps,\epsM,\bpsi,\bgam$ are stored as tensor components
$(a_{11},a_{22},a_{12})$ ($a_{12}$ tensorial). Their strain-like and stress-like Voigt vectors are
\begin{equation}
  \ve{\bm a}=\begin{bmatrix}a_{11}\\a_{22}\\2a_{12}\end{bmatrix},\qquad
  \vs{\bm a}=\begin{bmatrix}a_{11}\\a_{22}\\a_{12}\end{bmatrix},\qquad
  \bm a:\bm b=\vs{\bm a}^T\ve{\bm b},
  \label{eq:voigt}
\end{equation}
and stresses are $\vs{\bsig}=\Dm\ve{\eps}$. Throughout, $\bpsi$ denotes the symmetric
micro-deformation $\sym(\bpsi)$, and $\bm a^T\mathbb{X}\bm b:=\ve{\bm a}^T\mathbb{X}\ve{\bm b}$ for
any $3\times3$ stiffness $\mathbb{X}$. The micro-curvature is stored as the six-vector
\begin{equation}
  \bkap=[\kappa_{111},\kappa_{112},\kappa_{221},\kappa_{222},\kappa_{121},\kappa_{122}]^T,
  \qquad \kappa_{ijk}=\partial_k\psi_{ij}.
  \label{eq:kappa}
\end{equation}
\begin{delivers}
The conventions \eqref{eq:voigt}--\eqref{eq:kappa} used by every later step.
\end{delivers}

\subsection*{Step 0.3 -- Micromorphic strain energy}
With the macroscopic strain $\epsM=\sym(\nabla\bu)$, the relative strain
\begin{equation}
  \bgam=\epsM-\bpsi
  \label{eq:gammadef}
\end{equation}
and $\bkap=\nabla\bpsi$, the strain energy density is
\begin{equation}
  \boxed{\;W=\tfrac12\,\bpsi^T\C\,\bpsi+\tfrac12\,\bgam^T\Cb\,\bgam+\bpsi^T\Ch\,\bgam
  +\tfrac12\,\ell^2\,\bkap^T\GG\,\bkap\;}
  \label{eq:W}
\end{equation}
with the known homogenized stiffness $\C$ (Stage 2), and the unknown matrices of cubic symmetry
\begin{equation}
  \C=\begin{bmatrix}C_{11}&C_{12}&0\\C_{12}&C_{22}&0\\0&0&C_{33}\end{bmatrix},\quad
  \Cb=\begin{bmatrix}\bar C_{11}&\bar C_{12}&0\\\bar C_{12}&\bar C_{11}&0\\0&0&\bar C_{33}\end{bmatrix},\quad
  \Ch=\begin{bmatrix}\hat C_{11}&\hat C_{12}&0\\\hat C_{12}&\hat C_{11}&0\\0&0&\hat C_{33}\end{bmatrix},
  \label{eq:matrices}
\end{equation}
\begin{equation}
  \GG=\operatorname{diag}(G_{11},G_{22},G_{22},G_{11},G_{33},G_{33}),\qquad
  G_{11}=G_{22}=G_{33}=G\ \ (\code{Gmode='single'}).
  \label{eq:Gmat}
\end{equation}
The terms $C_{13},C_{23}$ of the computed $\C$ (zero by symmetry up to round-off) are discarded.
In components, using \eqref{eq:voigt},
\begin{align}
  2W={}&C_{11}\psi_{11}^2+2C_{12}\psi_{11}\psi_{22}+C_{22}\psi_{22}^2+4C_{33}\psi_{12}^2
  +\bar C_{11}(\gamma_{11}^2+\gamma_{22}^2)+2\bar C_{12}\gamma_{11}\gamma_{22}+4\bar C_{33}\gamma_{12}^2
  \nonumber\\
  &+2\hat C_{11}(\psi_{11}\gamma_{11}+\psi_{22}\gamma_{22})
  +2\hat C_{12}(\psi_{11}\gamma_{22}+\psi_{22}\gamma_{11})+8\hat C_{33}\psi_{12}\gamma_{12}
  +\ell^2G\,|\bkap|^2 .
  \label{eq:Wcomp}
\end{align}

\paragraph{Conjugate stresses.} Differentiating \eqref{eq:W} with $\bpsi$ and $\bgam$ treated as
independent (Voigt stress-like vectors):
\begin{equation}
  \bm s_{(\mathrm{sym})}=\frac{\partial W}{\partial\ve{\bpsi}}=\C\ve{\bpsi}+\Ch\ve{\bgam},\qquad
  \bsig_{(\gamma)}=\frac{\partial W}{\partial\ve{\bgam}}=\Cb\ve{\bgam}+\Ch\ve{\bpsi},\qquad
  \bmu=\frac{\partial W}{\partial\bkap}=\ell^2\GG\bkap .
  \label{eq:stresses}
\end{equation}
Since $\bgam=\epsM-\bpsi$, $\bsig_{(\gamma)}=\partial W/\partial\ve{\epsM}|_{\bpsi}$ is the stress
conjugate to the macroscopic strain.

\paragraph{Common eigenvectors.} Every symmetric matrix $\mathbb X$ with the structure
\eqref{eq:matrices} ($X_{11}=X_{22}$, $X_{13}=X_{23}=0$) is diagonalized by
\begin{equation}
\begin{gathered}
  \bm Q=[\bm v_1\ \bm v_2\ \bm v_3],\qquad
  \bm v_1=\tfrac{1}{\sqrt2}(1,1,0)^T,\quad \bm v_2=\tfrac{1}{\sqrt2}(1,-1,0)^T,\quad \bm v_3=(0,0,1)^T,\\
  \bm Q^T\mathbb X\bm Q=\operatorname{diag}(x_1,x_2,x_3),
\end{gathered}
  \label{eq:modes}
\end{equation}
with mode values $x_1=X_{11}+X_{12}$, $x_2=X_{11}-X_{12}$, $x_3=X_{33}$. For $\C$ we use
$a_C=\tfrac12(C_{11}+C_{22})$: $c_1=a_C+C_{12}$, $c_2=a_C-C_{12}$, $c_3=C_{33}$; likewise
$\bar c_i$ for $\Cb$ and $\hat c_i$ for $\Ch$.

\begin{proposition}[Positive definiteness of the local energy]
\label{prop:pd}
Let $\C\succ0$. The local part of \eqref{eq:W},
$\tfrac12[\ve{\bpsi};\ve{\bgam}]^T\B[\ve{\bpsi};\ve{\bgam}]$ with
\begin{equation}
  \B=\begin{bmatrix}\C&\Ch\\\Ch&\Cb\end{bmatrix},
  \label{eq:block}
\end{equation}
is positive definite if and only if, for every mode $i=1,2,3$,
\begin{equation}
  c_i\,\bar c_i-\hat c_i^{\,2}>0 .
  \label{eq:pdmode}
\end{equation}
\end{proposition}
\begin{proof}
With $\bm Q$ of \eqref{eq:modes}, $\operatorname{diag}(\bm Q,\bm Q)^T\B\operatorname{diag}(\bm Q,\bm Q)$
decouples into the three $2\times2$ blocks
$\big[\begin{smallmatrix}c_i&\hat c_i\\\hat c_i&\bar c_i\end{smallmatrix}\big]$, each positive
definite iff $c_i>0$ and $c_i\bar c_i-\hat c_i^2>0$ (Sylvester).
\end{proof}
\noindent In particular \eqref{eq:pdmode} implies $\bar c_i>0$, i.e.\ $\Cb\succ0$.
\begin{delivers}
Energy \eqref{eq:W}, conjugate stresses \eqref{eq:stresses}, unknowns
$\bth=[\bar C_{11},\bar C_{12},\bar C_{33},\hat C_{11},\hat C_{12},\hat C_{33},G]$ (7; 9 with
\code{Gmode='three'}), and the mode form \eqref{eq:pdmode} of positive definiteness.
\end{delivers}

% ==================================================================
\section{Stage 1 -- Load cases and kinematic boundary conditions}
\label{sec:stage1}
% ==================================================================

\subsection*{Step 1.1 -- From the raw load vector to $\bH$ and $\bG$}
\implementedin{\code{assembleHG}}
A load case is a raw vector $\hat{\bm h}=[\hat h_1,\dots,\hat h_9]$ ordered as
\[
  [H_{11},H_{22},H_{12},G_{1,11},G_{1,22},G_{1,12},G_{2,11},G_{2,22},G_{2,12}] .
\]
With amplitude $\varepsilon_0$ and gradient scale $g$ ($\varepsilon_0=0.01$, $g=2$):
\begin{equation}
\begin{aligned}
  &H_{11}=\varepsilon_0\hat h_1,\quad H_{22}=\varepsilon_0\hat h_2,\quad H_{12}=H_{21}=\varepsilon_0\hat h_3,\\
  &G_{111}=g\varepsilon_0\hat h_4,\ G_{122}=g\varepsilon_0\hat h_5,\ G_{112}=G_{121}=g\varepsilon_0\hat h_6,\\
  &G_{211}=g\varepsilon_0\hat h_7,\ G_{222}=g\varepsilon_0\hat h_8,\ G_{212}=G_{221}=g\varepsilon_0\hat h_9 .
\end{aligned}
\label{eq:HG}
\end{equation}

\subsection*{Step 1.2 -- Kinematic uniform boundary conditions (KUBC)}
All boundary nodes of every mesh receive
\begin{equation}
  u_i(\bX)=H_{ij}x_j+\tfrac12 G_{ijk}x_jx_k,\qquad \bX\in\partial\Omega,\quad \bx=\bX-\bX_c,
  \label{eq:kubc}
\end{equation}
whose strain is $\sym(\nabla\bu)_{ij}=\sym(\bH)_{ij}+\tfrac12(G_{ijk}+G_{jik})x_k$.

\subsection*{Step 1.3 -- Load-case library}
\implementedin{\code{buildCaseLibrary} (\code{cfg.cases.library})}
All library vectors have unit norm. \code{H3G16} contains 19 cases:

\begin{center}\small
\begin{tabular}{@{}cll@{}}
\toprule
\# & case & non-zero raw components\\
\midrule
1--3 & $H_{11}$, $H_{22}$, $H_{12}$ & $\hat h_1=1$; $\hat h_2=1$; $\hat h_3=1$\\
4--9 & single $G$ components & $\hat h_a=1$, $a=4,\dots,9$\\
10, 11 & bending / shear gradient in $X$ & $(\hat h_6,\hat h_7)=(1,\mp1)/\sqrt2$\\
12, 13 & bending / shear gradient in $Y$ & $(\hat h_9,\hat h_5)=(1,\mp1)/\sqrt2$\\
14, 15 & $G_{1,11}\pm G_{2,12}$ & $(\hat h_4,\hat h_9)=(1,\pm1)/\sqrt2$\\
16, 17 & $G_{1,12}\pm G_{2,22}$ & $(\hat h_6,\hat h_8)=(1,\pm1)/\sqrt2$\\
18 & $G_{1,11}+G_{2,22}$ & $(\hat h_4,\hat h_8)=(1,1)/\sqrt2$\\
19 & all six $G$ equal & $\hat h_4=\dots=\hat h_9=1/\sqrt6$\\
\bottomrule
\end{tabular}
\end{center}
\code{G16} omits cases 1--3, \code{G6} keeps the single $G$ components, \code{G21} contains the six
singles and all 15 pairs $(\hat{\bm e}_a+\hat{\bm e}_b)/\sqrt2$.
\begin{delivers}
For every load case $k\in\Kset$: $\bH^{(k)}$, $\bG^{(k)}$ and the boundary data \eqref{eq:kubc}.
\end{delivers}

% ==================================================================
\section{Stage 2 -- Homogenized stiffness $\C$ by periodic homogenization}
\label{sec:stage2}
% ==================================================================
\implementedin{\code{homogenize2D\_PBC.m} (called once, independent of $N$)}

\subsection*{Step 2.1 -- Periodic displacement field}
A single unit cell $Y=[0,L_c]^2$ with one centred inclusion is meshed with the same T6 elements as
in Stage 3 (periodic node pattern on opposite edges). For a prescribed macroscopic strain
$\bar\eps$ the displacement is
\begin{equation}
\begin{gathered}
  \bu(\bX)=\bar\eps\,\bX+\tilde\bu(\bX),\qquad \tilde\bu\ \text{periodic}\\
  \Longleftrightarrow\quad
  \bu(\bX+\bm L_\alpha)-\bu(\bX)=\bar\eps\,\bm L_\alpha,\qquad \bm L_1=(L_c,0),\ \bm L_2=(0,L_c).
\end{gathered}
  \label{eq:pbc}
\end{equation}

\subsection*{Step 2.2 -- Master/slave reduction}
Right- and top-edge nodes are slaved to their left/bottom partners, corners $(L_c,0)$, $(0,L_c)$,
$(L_c,L_c)$ to the origin corner, which is pinned to remove the rigid-body motion:
\begin{equation}
  \bu=\bm T\bu_m+\bu_p,\qquad
  \bu_p\big|_{\text{slave}}=\bar\eps\,\bm L_\alpha\ \text{(offsets)},
\end{equation}
where $\bm T$ maps the master degrees of freedom $\bu_m$ (interior, left and bottom edges) to all
degrees of freedom. The reduced equilibrium equations are
\begin{equation}
  \bm T^T\bK\bm T\,\bu_m=-\bm T^T\bK\,\bu_p .
  \label{eq:pbcsolve}
\end{equation}

\subsection*{Step 2.3 -- Three unit-strain cases and $\C$}
For $\ve{\bar\eps}\in\{(1,0,0),(0,1,0),(0,0,1)\}$ the volume averages
\begin{equation}
  \langle\vs{\bsig}\rangle=\frac1{V_c}\int_Y\vs{\bsig}\dd V,\qquad
  \langle\ve{\eps}\rangle=\frac1{V_c}\int_Y\ve{\eps}\dd V
\end{equation}
are computed by Gauss quadrature; the $k$-th column of $\C$ is
$\big(\langle\sigma_{11}\rangle,\langle\sigma_{22}\rangle,\langle\sigma_{12}\rangle\big)^T$, so that
\begin{equation}
  \langle\vs{\bsig}\rangle=\C\,\langle\ve{\eps}\rangle .
\end{equation}
The recovered strain is checked, $\|\langle\ve{\eps}\rangle-\ve{\bar\eps}\|<10^{-8}$. The
out-of-plane stress $\sigma_{33}=\nu(\sigma_{11}+\sigma_{22})$ of each phase yields the row
$\bm C_{s33}$ (not used in the identification).
\begin{delivers}
$\C=\C_{\mathrm{hom}}\in\mathbb R^{3\times3}$, with $\lambda_{\max}(\C)$ used later for scaling and
margins.
\end{delivers}

% ==================================================================
\section{Stage 3 -- Finite-element discretization}
\label{sec:stage3}
% ==================================================================

\subsection*{Step 3.1 -- Geometry and meshing}
\implementedin{\code{buildInclusionCenters}, \code{buildRUCMesh}, \code{buildPlainMesh}}
For a shift $\bm\zeta=(\zeta_1,\zeta_2)$ the inclusion centres are
\begin{equation}
  \bX^{(ij)}=\big((i+\tfrac12)\ell+\zeta_1,\ (j+\tfrac12)\ell+\zeta_2\big),\qquad i,j=-1,\dots,N,
  \label{eq:centres}
\end{equation}
kept if their bounding box intersects $\Omega$. The geometry is
$\Omega\cap\big(\Omega\cup\bigcup_{ij}B_R(\bX^{(ij)})\big)$ (set formula \code{(R1+C1+...)*R1}),
so inclusions crossing $\partial\Omega$ are clipped. The mesh has element size $h=0.05\,\ell$
(heterogeneous) or $h=0.02\,L$ (plain domain of Stage 5). An element belongs to the inclusion
phase iff its corner-node centroid lies inside one of the circles.

\subsection*{Step 3.2 -- Six-node triangle}
\implementedin{\code{t6ShapeDeriv}, \code{gauss3}, \code{buildElemData}}
With area coordinates $(L_1,L_2,L_3=1-L_1-L_2)$,
\begin{equation}
  \bm N=\big[L_1(2L_1-1),\ L_2(2L_2-1),\ L_3(2L_3-1),\ 4L_1L_2,\ 4L_2L_3,\ 4L_3L_1\big]^T,
  \label{eq:shape}
\end{equation}
Jacobian $\bm J=\begin{bmatrix}\partial_{L_1}\bm N^T\bm x_e&\partial_{L_2}\bm N^T\bm x_e\\
\partial_{L_1}\bm N^T\bm y_e&\partial_{L_2}\bm N^T\bm y_e\end{bmatrix}$ and
$[\partial_X\bm N,\ \partial_Y\bm N]=[\partial_{L_1}\bm N,\ \partial_{L_2}\bm N]\,\bm J^{-1}$.
The three-point rule uses $(L_1,L_2)_g\in\{(\tfrac23,\tfrac16),(\tfrac16,\tfrac23),(\tfrac16,\tfrac16)\}$,
$w_g=\tfrac16$. The strain matrix is
\begin{equation}
  \bm B=\begin{bmatrix}
  \partial_XN_1&0&\cdots&\partial_XN_6&0\\
  0&\partial_YN_1&\cdots&0&\partial_YN_6\\
  \partial_YN_1&\partial_XN_1&\cdots&\partial_YN_6&\partial_XN_6\end{bmatrix},\qquad
  \ve{\eps}=\bm B\bu_e .
  \label{eq:B}
\end{equation}

\subsection*{Step 3.3 -- Stiffness, KUBC solve and energy}
\implementedin{\code{elementStiffnessLE}, \code{assembleGlobalK\_LE}, \code{solveKUBC}}
\begin{equation}
  \bK_e=\sum_{g=1}^3w_g\det\bm J_g\,t\,\bm B_g^T\Dm_e\bm B_g,\qquad
  \bK=\operatorname*{\mathsf{A}}_e\bK_e .
  \label{eq:K}
\end{equation}
With the boundary degrees of freedom $d$ prescribed by \eqref{eq:kubc} and the free ones $f$,
\begin{equation}
  \bK_{ff}\,\bu_f=-\bK_{fd}\,\bu_d ,
  \label{eq:solve}
\end{equation}
solved for all load cases at once (one right-hand side per case). The strain energy is
\begin{equation}
  2U=\bu^T\bK\bu .
  \label{eq:energy}
\end{equation}
\begin{delivers}
For each mesh and load case: nodal displacements $\bu^{(k)}$ and energy $2U^{(k)}$.
\end{delivers}

% ==================================================================
\section{Stage 4 -- Evaluation grid, quadrature and regions}
\label{sec:stage4}
% ==================================================================

\subsection*{Step 4.1 -- Grid and point location}
\implementedin{\code{locateGridPoints}, \code{gridStrainOperator}}
All fields are evaluated on the regular grid $\bX_p=(x_a,y_b)$, $x_a=\tfrac{a-1}{n_g-1}L_x$,
$y_b=\tfrac{b-1}{n_g-1}L_y$, spacing $h_g$, with
\begin{equation}
  n_g(N)=\max\big(n_{g,\min},\ n_{\mathrm{cell}}N+1\big),\qquad n_{g,\min}=121,\ n_{\mathrm{cell}}=40,
\end{equation}
so that every unit cell is resolved by at least $n_{\mathrm{cell}}$ grid intervals for every $N$
(a fixed grid would make the quadrature error of the interface jumps grow with $N$). Each point is located in the straight-sided
(corner-node) triangle of an element through its barycentric coordinates
\begin{equation}
  \begin{bmatrix}L_1\\L_2\end{bmatrix}=
  \begin{bmatrix}x_1-x_3&x_2-x_3\\y_1-y_3&y_2-y_3\end{bmatrix}^{-1}
  \begin{bmatrix}X-x_3\\Y-y_3\end{bmatrix},\qquad L_3=1-L_1-L_2,
\end{equation}
the element with the largest $\min(L_1,L_2,L_3)$ being chosen (points on shared edges, interfaces).
The strain at $\bX_p$ is the element strain there, in tensor components,
\begin{equation}
  \varepsilon_{11}=\textstyle\sum_a\partial_XN_a u_{1a},\quad
  \varepsilon_{22}=\sum_a\partial_YN_a u_{2a},\quad
  \varepsilon_{12}=\tfrac12\sum_a(\partial_YN_a u_{1a}+\partial_XN_a u_{2a}),
\end{equation}
assembled into a sparse operator $\bm B_{\mathrm{grid}}$ with $\eps(\bX_p)=\bm B_{\mathrm{grid}}\bu$.
Jumps across material interfaces are retained (no nodal smoothing).

\subsection*{Step 4.2 -- Region weights}
\implementedin{\code{regionWeights}, \code{cellW}}
For a box $\mathcal R=[a_x,b_x]\times[a_y,b_y]$, each grid point receives the length of its grid
cell inside the box,
\begin{equation}
  w_a^{(x)}=\big|\,[x_a^-,x_a^+]\cap[a_x,b_x]\,\big|,\qquad
  x_a^\pm=\tfrac12(x_a+x_{a\pm1}),\ x_1^-=x_1,\ x_{n_g}^+=x_{n_g},
\end{equation}
likewise $w_b^{(y)}$, and
\begin{equation}
  \int_{\mathcal R}f\dd V\;\approx\;t\sum_p W_p f(\bX_p),\qquad W_p=w_a^{(x)}w_b^{(y)} .
  \label{eq:quad}
\end{equation}
For $\mathcal R=\Omega$ this is the trapezoidal rule; for fractional box edges the box area is
reproduced exactly.

\subsection*{Step 4.3 -- Full domain and interior region}
Every integral is evaluated over two regions:
\begin{equation}
  \mathcal R_{\mathrm{full}}=\Omega,\qquad
  \mathcal R_{\mathrm{int}}(N)=[a,L_x-a]\times[a,L_y-a],\quad a=\beta\,\ell ,
  \label{eq:regions}
\end{equation}
with band width $\beta$ (\code{cfg.interiorBand}, default $0.5$ cell). The interior excludes the
KUBC boundary layer; it is non-empty iff $2a<\min(L_x,L_y)$ (for $\beta=0.5$: $N\ge2$). Region
volume and root-mean-square moment arm:
\begin{equation}
  V_{\mathcal R}=(L_x-2a)(L_y-2a)\,t,\qquad
  L_{r,\mathcal R}=\frac{\max(L_x-2a,\ L_y-2a)}{\sqrt{12}}
  \label{eq:VL}
\end{equation}
($a=0$ for the full domain; $L/\sqrt{12}$ is the rms of $x$ over $[-L/2,L/2]$).
\begin{delivers}
Operator $\bm B_{\mathrm{grid}}$, weights $W_p$ of both regions, $V_{\mathcal R}(N)$,
$L_{r,\mathcal R}(N)$ and the validity flag of the interior region.
\end{delivers}

% ==================================================================
\section{Stage 5 -- Macroscopic reference field}
\label{sec:stage5}
% ==================================================================
\subsection*{Step 5.1 -- Homogeneous KUBC solve}
The plain domain with $\Dm\equiv\C$ is solved with \eqref{eq:kubc} for all load cases:
\begin{equation}
  \epsM^{(k)}(\bX_p)=\bm B_{\mathrm{grid}}\bu_{\mathrm{hom}}^{(k)},\qquad
  2U^{(k)}_{\mathrm{hom}}=\bu_{\mathrm{hom}}^{(k)T}\bK_{\C}\bu_{\mathrm{hom}}^{(k)} .
  \label{eq:stage1}
\end{equation}
Neither depends on $N$. For the interior region, $2U^{(k)}_{\mathrm{hom},\mathrm{int}}
=\bu_{\mathrm{hom}}^{(k)T}\bK_{r}\bu_{\mathrm{hom}}^{(k)}$ with the region stiffness $\bK_r$ of
Step 6.3 (depends on $N$ through $a$).
\begin{delivers}
$\epsM^{(k)}$ on the grid, $2U^{(k)}_{\mathrm{hom}}$ for both regions.
\end{delivers}

% ==================================================================
\section{Stage 6 -- Translation averaging and FE stress integrals}
\label{sec:stage6}
% ==================================================================

\subsection*{Step 6.1 -- Shift set over one full period}
\implementedin{\code{buildTranslationList} (mode 5)}
\begin{equation}
  \Zset=\big\{(s_a\ell,\ s_b\ell)\ :\ s_a=\tfrac{a+1/2}{n}-\tfrac12,\ a,b=0,\dots,n-1\big\},
  \qquad|\Zset|=n^2,
  \label{eq:shifts}
\end{equation}
($n=8$ by default). $\Zset$ is uniform over one cell and symmetric, $\bm\zeta\in\Zset\Rightarrow
-\bm\zeta\in\Zset$.

\subsection*{Step 6.2 -- One heterogeneous solve per shift}
For every $\bm\zeta\in\Zset$: mesh (Step 3.1), assemble $\bK^{(\bm\zeta)}$, solve \eqref{eq:solve}
for all load cases, and accumulate
\begin{equation}
  \eps^{(k,\bm\zeta)}(\bX_p)=\bm B^{(\bm\zeta)}_{\mathrm{grid}}\bu^{(k,\bm\zeta)},\qquad
  2U^{(k,\bm\zeta)}=\bu^{(k,\bm\zeta)T}\bK^{(\bm\zeta)}\bu^{(k,\bm\zeta)},
\end{equation}
and the region integrals of Step 6.3.

\subsection*{Step 6.3 -- Exact region integrals of the FE stress}
\implementedin{\code{regionOps}}
For a box $\mathcal R$, only Gauss points $\bX_g\in\mathcal R$ contribute. With
$\omega_g=w_g\det\bm J_g\,t$ and $\bx_g=\bX_g-\bX_c$,
\begin{align}
  \int_{\mathcal R}\vs{\bsig}\dd V&=\bm Q_s\bu, &
  \bm Q_s&=\operatorname*{\mathsf{A}}_e\sum_{g:\,\bX_g\in\mathcal R}\omega_g\,\Dm_e\bm B_g,
  \label{eq:Qs}\\
  \int_{\mathcal R}\sigma_{ij}x_k\dd V&=\bm Q_m\bu, &
  \bm Q_m&=\operatorname*{\mathsf{A}}_e\sum_{g:\,\bX_g\in\mathcal R}\omega_g
  \begin{bmatrix}(\Dm_e\bm B_g)_{1\cdot}\,x_{1g}\\(\Dm_e\bm B_g)_{1\cdot}\,x_{2g}\\
  (\Dm_e\bm B_g)_{2\cdot}\,x_{1g}\\(\Dm_e\bm B_g)_{2\cdot}\,x_{2g}\\
  (\Dm_e\bm B_g)_{3\cdot}\,x_{1g}\\(\Dm_e\bm B_g)_{3\cdot}\,x_{2g}\end{bmatrix},
  \label{eq:Qm}\\
  \int_{\mathcal R}\bsig:\eps\dd V&=\bu^T\bK_r\bu, &
  \bK_r&=\operatorname*{\mathsf{A}}_e\sum_{g:\,\bX_g\in\mathcal R}\omega_g\,\bm B_g^T\Dm_e\bm B_g ,
  \label{eq:Kr}
\end{align}
the moment vector being ordered
$[\sigma_{11}x_1,\sigma_{11}x_2,\sigma_{22}x_1,\sigma_{22}x_2,\sigma_{12}x_1,\sigma_{12}x_2]$ and
$x_{kg}=\bm N_g^T\bm X_{k,e}-X_{c,k}$. For $\mathcal R=\Omega$, $\bK_r=\bK$.

\subsection*{Step 6.4 -- Averages over the shifts}
For every load case and region,
\begin{align}
  \bpsi^{(k)}(\bX_p)&=\frac1{|\Zset|}\sum_{\bm\zeta}\eps^{(k,\bm\zeta)}(\bX_p),
  \label{eq:psi}\\
  2U^{(k)}_{\mathrm{het},\mathcal R}&=\frac1{|\Zset|}\sum_{\bm\zeta}\bu^{(k,\bm\zeta)T}\bK_r^{(\bm\zeta)}\bu^{(k,\bm\zeta)}
  \quad(\text{and its standard deviation}),\\
  \bm S^{(k)}_{\mathrm{FE},\mathcal R}&=\frac1{|\Zset|}\sum_{\bm\zeta}\bm Q_s^{(\bm\zeta)}\bu^{(k,\bm\zeta)}\in\mathbb R^3,\qquad
  \bm M^{(k)}_{\mathrm{FE},\mathcal R}=\frac1{|\Zset|}\sum_{\bm\zeta}\bm Q_m^{(\bm\zeta)}\bu^{(k,\bm\zeta)}\in\mathbb R^6 .
  \label{eq:SMfe}
\end{align}
Load cases are never averaged together.
\begin{delivers}
$\bpsi^{(k)}$ on the grid; $2U^{(k)}_{\mathrm{het}}$, $\bm S^{(k)}_{\mathrm{FE}}=\int\sigma_{ij}$ and
$\bm M^{(k)}_{\mathrm{FE}}=\int\sigma_{ij}x_k$ for both regions.
\end{delivers}

% ==================================================================
\section{Stage 7 -- Relative strain and micro-curvature}
\label{sec:stage7}
% ==================================================================

\subsection*{Step 7.1 -- Relative strain}
\begin{equation}
  \bgam^{(k)}(\bX_p)=\epsM^{(k)}(\bX_p)-\bpsi^{(k)}(\bX_p) .
\end{equation}

\subsection*{Step 7.2 -- Micro-curvature: exact macroscopic part plus smoothed deviation (default)}
\implementedin{\code{gridStrainOperator} (Stage-1 gradients), \code{localGradient}, \code{kernelMats} (\code{cfg.kappaMethod='local'})}
Since $\bpsi=\epsM-\bgam$,
\begin{equation}
  \boxed{\;\bkap=\nabla\bpsi=\nabla\epsM-\nabla\bgam\;}
  \label{eq:kapsplit}
\end{equation}
The macroscopic part $\nabla\epsM$ is taken exactly from the Stage-1 FE solution (element-wise
strain gradients of the homogeneous, finely meshed field: for straight T6 elements
$\partial\eps/\partial\bX=[\eps_1-\eps_3,\ \eps_2-\eps_3]\,\bm A^{-1}$ with corner strains $\eps_m$
and chord Jacobian $\bm A$); it contains no ripple and no interfaces and is independent of $N$.
Only the small deviation $\bgam$ is differentiated with the moving least-squares operator below.
At every grid point $\bX_0$, each component $f$ of $\bgam$ is fitted locally by an affine function,
\begin{equation}
  \min_{a,\bm b}\ \sum_p w_p(\bX_0)\,\big(f_p-a-\bm b\cdot(\bX_p-\bX_0)\big)^2,\qquad
  w_p=g(X_p-X_0)\,g(Y_p-Y_0),\quad
  g(d)=e^{-d^2/(2\sigma^2)}\,\mathbf 1_{|d|\le3\sigma},
  \label{eq:mls}
\end{equation}
with $\sigma=w_\sigma\ell$ ($w_\sigma=$\code{kappaLocalWidth}$=0.1$), raised to $1.5h_g$ if smaller.
The gradient is $\nabla f(\bX_0)\approx\bm b$. The normal equations are
\begin{equation}
  \underbrace{\begin{bmatrix}m_{11}&m_{12}&m_{13}\\m_{12}&m_{22}&m_{23}\\m_{13}&m_{23}&m_{33}\end{bmatrix}}_{\bm M(\bX_0)}
  \begin{bmatrix}a\\b_1\\b_2\end{bmatrix}=
  \begin{bmatrix}r_1\\r_2\\r_3\end{bmatrix},\qquad
  \begin{aligned}
  &m_{11}=\textstyle\sum w,\ m_{12}=\sum w\,\delta x,\ m_{13}=\sum w\,\delta y,\\
  &m_{22}=\textstyle\sum w\,\delta x^2,\ m_{23}=\sum w\,\delta x\,\delta y,\ m_{33}=\sum w\,\delta y^2,\\
  &r_1=\textstyle\sum wf,\ r_2=\sum wf\,\delta x,\ r_3=\sum wf\,\delta y,
  \end{aligned}
\end{equation}
$\delta x=X_p-X_0$, $\delta y=Y_p-Y_0$. Because $w$ is separable, all sums are matrix products with
the one-dimensional kernel matrices $\bm K^{(q)}_{ij}=g(x_j-x_i)\,(x_j-x_i)^q$, $q=0,1,2$ (on the
grid array $\bm F$ of $f$):
\begin{equation}
  r_1=\bm K_y^{(0)}\bm F\bm K_x^{(0)T},\quad r_2=\bm K_y^{(0)}\bm F\bm K_x^{(1)T},\quad
  r_3=\bm K_y^{(1)}\bm F\bm K_x^{(0)T},
\end{equation}
and the $m_{ij}$ likewise with $\bm F=\bm 1$. With the cofactors
$c_{11}=m_{22}m_{33}-m_{23}^2$, $c_{12}=-(m_{12}m_{33}-m_{13}m_{23})$, $c_{22}=m_{11}m_{33}-m_{13}^2$,
$c_{23}=-(m_{11}m_{23}-m_{12}m_{13})$, $c_{13}=m_{12}m_{23}-m_{13}m_{22}$,
$c_{33}=m_{11}m_{22}-m_{12}^2$ and $\det\bm M=m_{11}c_{11}+m_{12}c_{12}+m_{13}c_{13}$,
\begin{equation}
  b_1=\frac{c_{12}r_1+c_{22}r_2+c_{23}r_3}{\det\bm M},\qquad
  b_2=\frac{c_{13}r_1+c_{23}r_2+c_{33}r_3}{\det\bm M} .
  \label{eq:mlssol}
\end{equation}
Applied to $\gamma_{11},\gamma_{22},\gamma_{12}$ this gives $\nabla\bgam$, and $\bkap$ follows from
\eqref{eq:kapsplit} in the order \eqref{eq:kappa}.

\begin{lemma}[Properties of the local gradient]
\label{lem:mls}
(i) For an affine field $f=\alpha+\bm\beta\cdot\bX$ the estimate is exact, $\bm b=\bm\beta$, at
every grid point including the boundary, since the fit is unbiased for affine data regardless of
the (one-sided) weight distribution. (ii) Because $\sigma\propto\ell$, a microstructural boundary
layer of width proportional to $\ell$ is resolved to the same relative accuracy for every $N$.
Macroscopic features (width independent of $\ell$) would be resolved differently at every $N$;
this is why the macroscopic part is excluded from the window in \eqref{eq:kapsplit}.
(iii) The ripple with period $\ell/n$ left by the finite shift set is damped by
$\exp(-2\pi^2\sigma^2n^2/\ell^2)$, which requires $w_\sigma\gtrsim0.6/n$
($e^{-12.6}$ for $n=8$, $w_\sigma=0.1$).
\end{lemma}

\subsection*{Step 7.3 -- Alternatives}
\code{'localPlain'}: the window applied to $\bpsi$ itself (resolves the macroscopic boundary
gradients differently at every $N$). \code{'polyfit'}: weighted least-squares fit of each component with the monomials
$\phi_{ab}=s^a\tau^b$, $a+b\le p$, $s=(X-X_c)/L_x$, $\tau=(Y-Y_c)/L_y$,
$\bm c=(\bm\Phi^T\bm W\bm\Phi)^{-1}\bm\Phi^T\bm W\bm f$, differentiated analytically (removes
boundary layers). \code{'element'}: element-wise strain gradients averaged over the shifts
(misses the interface jumps). \code{'gridFD'}: central differences (grid dependent).

\subsection*{Step 7.4 -- Curvature of the macroscopic field}
$\bkap_M=\nabla\epsM$ is the exact Stage-1 gradient of \eqref{eq:kapsplit}, used for the reference
integrals of Stage 8. As a check of the gradient operator, the code reports the variation of
$\int|\bkap_M|^2$ over $N$ (full domain), which must vanish up to quadrature
($\approx1\%$; with the window applied to $\epsM$ it reached $>100\%$).
\begin{delivers}
$\bgam^{(k)}$, $\bkap^{(k)}$ and $\bkap_M^{(k)}$ on the grid.
\end{delivers}

% ==================================================================
\section{Stage 8 -- Integral data}
\label{sec:stage8}
% ==================================================================
\implementedin{\code{dataGeneration}, Stage 4 of the code}
With \eqref{eq:quad} over region $\mathcal R$ (all tensor components):
\begin{align}
  &A_1=\!\int\!\psi_{11}^2,\ A_2=\!\int\!\psi_{11}\psi_{22},\ A_3=\!\int\!\psi_{22}^2,\ A_4=\!\int\!\psi_{12}^2,
  \nonumber\\
  &A_5=\!\int\!\gamma_{11}^2,\ A_6=\!\int\!\gamma_{11}\gamma_{22},\ A_7=\!\int\!\gamma_{22}^2,\ A_8=\!\int\!\gamma_{12}^2,
  \nonumber\\
  &A_9=\!\int\!\psi_{11}\gamma_{11},\ A_{10}=\!\int\!\psi_{11}\gamma_{22},\ A_{11}=\!\int\!\psi_{22}\gamma_{11},\
  A_{12}=\!\int\!\psi_{22}\gamma_{22},\ A_{13}=\!\int\!\psi_{12}\gamma_{12},
  \nonumber\\
  &A_{14}=\!\int\!\kappa_{111}^2,\ A_{15}=\!\int\!\kappa_{112}^2,\ A_{16}=\!\int\!\kappa_{122}^2,\
  A_{17}=\!\int\!\kappa_{222}^2,\ A_{18}=\!\int\!\kappa_{121}^2,\ A_{19}=\!\int\!\kappa_{221}^2,
  \label{eq:A}
\end{align}
\begin{align}
  &I_{1..3}=\!\int\!(\psi_{11},\psi_{22},\psi_{12}),\quad
  I_{4..6}=\!\int\!(\gamma_{11},\gamma_{22},\gamma_{12}),\quad
  I_{7..12}=\!\int\!(\kappa_{111},\kappa_{112},\kappa_{221},\kappa_{222},\kappa_{121},\kappa_{122}),
  \nonumber\\
  &I_{13..18}=\!\int\!(\psi_{11}x_1,\psi_{11}x_2,\psi_{22}x_1,\psi_{22}x_2,\psi_{12}x_1,\psi_{12}x_2),
  \nonumber\\
  &I_{19..24}=\!\int\!(\gamma_{11}x_1,\gamma_{11}x_2,\gamma_{22}x_1,\gamma_{22}x_2,\gamma_{12}x_1,\gamma_{12}x_2),
  \label{eq:I}
\end{align}
with $x_k$ measured from the centre of $\Omega$ for both regions. The macroscopic references
$A^M_j$ are the integrals of $A_1$--$A_4$ with $\bpsi\to\epsM$ and of $A_{14}$--$A_{19}$ with
$\bkap\to\bkap_M$ ($A^M_5\dots A^M_{13}$ are not defined). As a quadrature check,
$\int_\Omega\epsM^T\C\epsM\dd V$ (grid) is compared with $2U_{\mathrm{hom}}$ (FE).
\begin{delivers}
For every region, $N$ and load case: $A_1\dots A_{19}$, $I_1\dots I_{24}$, $A^M_j$,
$2U_{\mathrm{het}}$, $2U_{\mathrm{hom}}$, $\bm S_{\mathrm{FE}}$, $\bm M_{\mathrm{FE}}$,
$V_{\mathcal R}$, $L_{r,\mathcal R}$, $\ell^2$ (file \code{DataGeneration.mat}).
\end{delivers}

% ==================================================================
\section{Stage 9 -- Smoothing over $N$}
\label{sec:stage9}
% ==================================================================
\implementedin{\code{smoothData}, \code{fitN}, \code{nnlsScaled}}

\subsection*{Step 9.1 -- Region and admissible $N$}
The chosen region's data are mapped onto the standard names; only
\[
  N\in\Nset_s=\{N\ge N_{\min},\ \text{region non-empty}\}
\]
are used ($N_{\min}=1$ by default; at least three values are required).

\subsection*{Step 9.2 -- Densities}
To remove the purely geometric $N$-dependence of the interior region, every quantity is fitted as a
density and multiplied back afterwards:
\begin{equation}
  \tilde A_j=\frac{A_j}{V_{\mathcal R}}\ (j\le13),\quad
  \tilde A_j=\frac{\ell^2A_j}{V_{\mathcal R}}\ (j\ge14),\quad
  \tilde I_j=\frac{I_j}{V_{\mathcal R}}\ (j\le12),\quad
  \tilde I_j=\frac{I_j}{V_{\mathcal R}L_{r,\mathcal R}}\ (j\ge13),
\end{equation}
$\tilde{\bm S}=\bm S_{\mathrm{FE}}/V_{\mathcal R}$, $\tilde{\bm M}=\bm M_{\mathrm{FE}}/(V_{\mathcal R}L_{r,\mathcal R})$.
The curvature terms are fitted as $\ell^2A_j$ because $A_j$ itself may grow with $N$ (a boundary
layer of width $\sim\ell$ and gradient $\sim\Delta\psi/\ell$ gives $\int\kappa^2\sim N$), while the
energy contribution $\ell^2A_j$ decreases (bulk $\sim N^{-2}$, boundary layer $\sim N^{-1}$).

\subsection*{Step 9.3 -- Shape-constrained fits}
Every density $y(N)$ is represented by
\begin{equation}
  \hat y(N)=\sum_kc_kN^{-p_k},
  \label{eq:fit}
\end{equation}
fitted by weighted relative least squares,
\begin{equation}
  \min_{\bm c}\sum_{N\in\Nset_s}\Big(\frac{\hat y(N)-y(N)}{\max(|y(N)|,\,10^{-3}\max|y|)}\Big)^2,
  \label{eq:wls}
\end{equation}
under the constraints of Table~\ref{tab:smooth}. Non-negativity is solved by \code{lsqnonneg} on
column- and row-balanced data ($\bm A\to\bm A\,\mathrm{diag}(1/\max_i|A_{ij}|)$,
$\bm y\to\bm y/\max|y|$); ``one sign'' fits $c\ge0$ to $y$ and to $-y$ and keeps the better. A fit
never has more terms than $|\Nset_s|-1$; if the exponent set is larger, the best-fitting subset of
that size is selected.

\begin{table}[h]
\centering\small
\begin{tabular}{@{}lll@{}}
\toprule
quantity (density) & model & property\\
\midrule
$A_5,A_7,A_8$ ($\gamma\gamma$) & $c_k\ge0$, $p\in\mathcal P_d$ & decreasing $\to0$\\
$A_6$, $A_9$--$A_{13}$ & one sign, $p\in\mathcal P_d$ & $|\cdot|$ decreasing $\to0$\\
$\ell^2A_{14}$--$\ell^2A_{19}$ & $c_k\ge0$, $p\in\mathcal P_d$ & decreasing $\to0$\\
$I_{19}$--$I_{24}$ ($\int\gamma x_k$) & one sign, $p\in\mathcal P_d$ & $|\cdot|$ decreasing $\to0$\\
$A_1$--$A_4$, $I_1$--$I_{18}$, $\bm S_{\mathrm{FE}}$, $\bm M_{\mathrm{FE}}$ & free, $p\in\mathcal P_f$ & smooth\\
$2U_{\mathrm{het}}/V-2U_{\mathrm{hom}}/V$ & one sign, $p\in\mathcal P_d$ & excess $\to0$\\
\bottomrule
\end{tabular}
\caption{Smoothing models, $\mathcal P_d=\{\tfrac12,1,\tfrac32,2\}$, $\mathcal P_f=\{0,1,2\}$.}
\label{tab:smooth}
\end{table}

\begin{lemma}[Monotonicity]
If $p_k>0$, $c_k\ge0$ and $\sum_kc_k>0$, then $\hat y$ of \eqref{eq:fit} is positive, strictly
decreasing and convex, and $\hat y(N)\to0$ for $N\to\infty$.
\end{lemma}
\begin{proof}
$\hat y'=-\sum_kc_kp_kN^{-p_k-1}<0$, $\hat y''=\sum_kc_kp_k(p_k+1)N^{-p_k-2}>0$.
\end{proof}

\subsection*{Step 9.4 -- Energy}
With the exact Stage-1 density $e_{\mathrm{hom}}(N)=2U_{\mathrm{hom},\mathcal R}/V_{\mathcal R}$,
the excess $2U_{\mathrm{het}}/V_{\mathcal R}-e_{\mathrm{hom}}$ is fitted (one sign) and
$2\hat U_{\mathrm{het}}=(e_{\mathrm{hom}}+\widehat{\text{excess}})V_{\mathcal R}$.
\begin{delivers}
Smoothed $A_j$, $I_j$, $\bm S_{\mathrm{FE}}$, $\bm M_{\mathrm{FE}}$, $2U_{\mathrm{het}}$ on $\Nset_s$
(file \code{SmoothedData.mat}); relative RMS of every fit.
\end{delivers}

% ==================================================================
\section{Stage 10 -- Identification equations}
\label{sec:stage10}
% ==================================================================
\implementedin{\code{caseRows}}
For every load case $k$ and $N$, three groups of linear equations in
\[
  \bth=[\bar C_{11},\bar C_{12},\bar C_{33},\hat C_{11},\hat C_{12},\hat C_{33},G]
\]
follow by integrating \eqref{eq:Wcomp} and \eqref{eq:stresses} over the region with the definitions
\eqref{eq:A}--\eqref{eq:I}.

\subsection*{Step 10.1 -- Energy (1 equation)}
\begin{equation}
  \boxed{
  \begin{aligned}
  &\bar C_{11}(A_5+A_7)+2\bar C_{12}A_6+4\bar C_{33}A_8
  +2\hat C_{11}(A_9+A_{12})+2\hat C_{12}(A_{10}+A_{11})+8\hat C_{33}A_{13}\\
  &\quad+\ell^2G\sum_{j=14}^{19}A_j
  \;=\;2U_{\mathrm{het}}-\big(C_{11}A_1+2C_{12}A_2+C_{22}A_3+4C_{33}A_4\big).
  \end{aligned}}
  \label{eq:eqE}
\end{equation}
With \code{Gmode='three'} the last term is
$\ell^2\big[G_{11}(A_{14}+A_{17})+G_{22}(A_{15}+A_{19})+G_{33}(A_{16}+A_{18})\big]$.

\subsection*{Step 10.2 -- Integrated stress (3 equations)}
The FE stress integral is matched with the integral of $\bsig_{(\gamma)}$, the stress conjugate to
$\epsM$:
\begin{equation}
  \boxed{
  \begin{aligned}
  \bar C_{11}I_4+\bar C_{12}I_5+\hat C_{11}I_1+\hat C_{12}I_2&=S_{\mathrm{FE},1},\\
  \bar C_{12}I_4+\bar C_{11}I_5+\hat C_{12}I_1+\hat C_{11}I_2&=S_{\mathrm{FE},2},\\
  2\bar C_{33}I_6+2\hat C_{33}I_3&=S_{\mathrm{FE},3}.
  \end{aligned}}
  \label{eq:eqS}
\end{equation}

\subsection*{Step 10.3 -- Integrated double stress (6 equations)}
For a quadratic virtual displacement $\delta u_i=\tfrac12\delta G_{ijk}x_jx_k$ the virtual strain is
$\delta\eps=\delta\bkap\,\bx$ with constant $\delta\kappa_{ijk}=\tfrac12(\delta G_{ijk}+\delta G_{jik})$,
and
\begin{equation}
  \int\bsig:\delta\eps\dd V=\sum_k\int\big(\sigma_{11}\delta\kappa_{11k}+\sigma_{22}\delta\kappa_{22k}
  +2\sigma_{12}\delta\kappa_{12k}\big)x_k\dd V=\delta\bkap\cdot\bm D\!\int\!\bm m(\bsig,\bx)\dd V,
\end{equation}
$\bm D=\operatorname{diag}(1,1,1,1,2,2)$, $\bm m$ the moment vector of \eqref{eq:Qm}. For a
second-gradient continuum with local stress $\bsig_{(\gamma)}$ and double stress $\bmu$ the internal
virtual work of the same field is $\delta\bkap\cdot\big(\bm D\int\bm m(\bsig_{(\gamma)},\bx)+\int\bmu\big)$.
Equating both for all $\delta\bkap$ gives (\code{doubleStressModel='local+mu'})
\begin{equation}
  \boxed{\ \bm D\!\int\!\bm m(\bsig_{(\gamma)},\bx)\dd V+\ell^2G\,[I_7,\dots,I_{12}]^T=\bm D\,\bm M_{\mathrm{FE}}\ }
  \label{eq:eqM}
\end{equation}
with the local part
\begin{equation}
  \bm D\!\int\!\bm m(\bsig_{(\gamma)},\bx)=
  \begin{bmatrix}
  \bar C_{11}I_{19}+\bar C_{12}I_{21}+\hat C_{11}I_{13}+\hat C_{12}I_{15}\\
  \bar C_{11}I_{20}+\bar C_{12}I_{22}+\hat C_{11}I_{14}+\hat C_{12}I_{16}\\
  \bar C_{12}I_{19}+\bar C_{11}I_{21}+\hat C_{12}I_{13}+\hat C_{11}I_{15}\\
  \bar C_{12}I_{20}+\bar C_{11}I_{22}+\hat C_{12}I_{14}+\hat C_{11}I_{16}\\
  4\bar C_{33}I_{23}+4\hat C_{33}I_{17}\\
  4\bar C_{33}I_{24}+4\hat C_{33}I_{18}
  \end{bmatrix}.
  \label{eq:eqMloc}
\end{equation}
(\code{'mu'} keeps only $\ell^2G[I_7..I_{12}]$ on the left.) With \code{Gmode='three'},
\[
  \int\bmu=\ell^2[G_{11}I_7,\ G_{22}I_8,\ G_{22}I_9,\ G_{11}I_{10},\ G_{33}I_{11},\ G_{33}I_{12}]^T .
\]

\subsection*{Step 10.4 -- Scaling, weights and informative groups}
Each group is divided by its own size and weighted:
\begin{equation}
  \text{energy: }\frac{\sqrt{w_E}}{2U_{\mathrm{het}}},\qquad
  \text{stress: }\frac{\sqrt{w_S}}{\|\bm S_{\mathrm{FE}}\|},\qquad
  \text{double stress: }\frac{\sqrt{w_D}}{\|\bm D\bm M_{\mathrm{FE}}\|}.
\end{equation}
A stress or double-stress group enters only if it carries information,
\begin{equation}
  \|\bm S_{\mathrm{FE}}\|\ge\tau\,\sigma_{\mathrm{ref}},\qquad
  \|\bm D\bm M_{\mathrm{FE}}\|\ge\tau\,\sigma_{\mathrm{ref}}L_{r,\mathcal R},\qquad
  \sigma_{\mathrm{ref}}=\sqrt{\frac{2U_{\mathrm{het}}\,\lambda_{\max}(\C)}{V_{\mathcal R}}}\;V_{\mathcal R},
  \label{eq:info}
\end{equation}
$\tau=10^{-2}$. Stacking all groups of all load cases gives the linear least-squares system
\begin{equation}
  \bP(N)\,\bth\approx\bb(N).
  \label{eq:system}
\end{equation}
\begin{delivers}
Rows $\bP(N)$ and right-hand sides $\bb(N)$ for every $N$, labelled by channel (energy, stress,
double stress).
\end{delivers}

% ==================================================================
\section{Stage 11 -- Solution: monotone in $N$, energy positive definite (default)}
\label{sec:stage11}
% ==================================================================
\implementedin{\code{solveMonotone}, \code{monoObjective}, \code{monoNonlcon} (\code{cfg.id.mode='monotone'})}

\subsection*{Step 11.1 -- Parameters as functions of $N$}
Every parameter is represented as
\begin{equation}
  \theta_j(N)=\sum_{k=1}^{n_b}c_{jk}N^{-e_k},\qquad c_{jk}\ge0,\qquad \bm e=(0,1,2),
  \label{eq:mono}
\end{equation}
hence positive and non-increasing in $N$ (Lemma of Stage 9 with $p=0$ allowed), with limit
$c_{j1}$ for $N\to\infty$. With $\bm\beta(N)=[N^{-e_1},\dots,N^{-e_{n_b}}]$ and the
parameter-major coefficient vector $\bm c$,
\begin{equation}
  \bth(N)=\bm T(N)\bm c,\qquad \bm T(N)=\bm I_{n_p}\otimes\bm\beta(N),
\end{equation}
and the stacked system over all $N\in\Nset_s$ is $\tilde\bP\bm c\approx\tilde\bb$,
$\tilde\bP=[\bP(N)\bm T(N)]_N$. Columns are scaled, $\bm c=\bm S\bm x$,
$S_{jj}=1/\max\big(\max_i|\tilde P_{ij}|,\ 10^{-6}\max_{ij}|\tilde P_{ij}|\big)$.

\subsection*{Step 11.2 -- Constraints at every $N$}
With the mode values of Step 0.3 as linear functions of $\bm c$,
$\hat c_i(N)=\bm u_i(N)^T\bm c$, $\bar c_i(N)=\bm v_i(N)^T\bm c$, margin
$m=\delta\,\lambda_{\max}(\C)$, $\delta=$\code{pdMargin}$=10^{-3}$:
\begin{align}
  &\hat c_i(N)\ge m &&(\Ch\succ0,\ \text{linear}), \label{eq:conlin}\\
  &\frac{\hat c_i(N)^2}{(1-\delta)\,c_i}-\bar c_i(N)\le0 &&(\B\succ0,\ \text{Prop.~\ref{prop:pd}},\ \text{convex}),
  \label{eq:connl}\\
  &\bm c\ge\bm0 &&(\text{positive, non-increasing}),
\end{align}
for $i=1,2,3$ and all $N\in\Nset_s$. \eqref{eq:connl} is the convex quadratic epigraph
$(\bm u_i^T\bm c)^2\le(1-\delta)c_i\,\bm v_i^T\bm c$ and implies $\bar c_i\ge\hat c_i^2/c_i>0$.
Without \code{energyPD}, \eqref{eq:connl} is replaced by the linear $\bar c_i(N)\ge m$.

\subsection*{Step 11.3 -- Convex optimization problem}
\begin{equation}
  \min_{\bm x\ge\bm0}\ J(\bm x)=\tfrac12\|\tilde\bP\bm S\bm x-\tilde\bb\|^2
  \quad\text{s.t.}\quad\eqref{eq:conlin},\ \eqref{eq:connl},
  \label{eq:monoopt}
\end{equation}
a convex quadratic program with convex quadratic constraints: every local minimum is global.
Gradients supplied to the solver:
\begin{equation}
  \nabla J=\bm S\tilde\bP^T(\tilde\bP\bm S\bm x-\tilde\bb),\qquad
  \nabla_{\bm x}\Big[\frac{(\bm u_i^T\bm S\bm x)^2}{(1-\delta)c_i}-\bm v_i^T\bm S\bm x\Big]
  =\frac{2\,\bm u_i^T\bm S\bm x}{(1-\delta)c_i}\bm S\bm u_i-\bm S\bm v_i .
\end{equation}
Start point (feasible): $\Ch=\C$, $\Cb=2\C$, $G=10^{-3}\lambda_{\max}(\C)$ on the lowest exponent.
Solver: \code{fmincon} (SQP) in MATLAB, \code{sqp} in Octave.

\subsection*{Step 11.4 -- Post-processing}
Coefficients below $10^{-10}\max|c_{jk}|$ are set to zero (bound $c\ge0$ at round-off); a parameter
with $c_{jk}=0$ for all $e_k>0$ is reported as constant in $N$. For every $N$:
$\bth(N)=\bm T(N)\bm c$; channel misfits
\begin{equation}
  \mathrm{RMS}_q(N)=\Big(\frac1{n_q}\sum_{i\in q,N}\frac{r_i^2}{w_q}\Big)^{1/2},\qquad
  \bm r=\tilde\bP\bm c-\tilde\bb ;
\end{equation}
and $\lambda_{\min}(\Cb)$, $\lambda_{\min}(\Ch)$, $\lambda_{\min}(\B)$, reported relative to
$\lambda_{\max}(\C)$.
\begin{delivers}
$\bth(N)$ for every $N$, the closed form $\theta_j(N)=c_{j1}+c_{j2}/N+c_{j3}/N^2$, channel misfits,
eigenvalue checks, constant-in-$N$ parameters (\code{IdentifiedCbarChatG.mat}, \code{.csv}).
\end{delivers}

% ==================================================================
\section{Stage 12 -- Solution: free per $N$ and global (comparison)}
\label{sec:stage12}
% ==================================================================
\implementedin{\code{solveBlock}, \code{pdConstraints}, \code{qpActiveSet}, \code{fitTrends}}

\subsection*{Step 12.1 -- Linearly constrained least squares}
For one $N$ (\code{'perN'}) or all $N$ together with one $\bth$ (\code{'global'}):
\begin{equation}
  \min_{\bth}\tfrac12\|\bP\bth-\bb\|^2\quad\text{s.t.}\quad \bm G_c\bth\ge\bm h,
  \label{eq:qp}
\end{equation}
with the rows $\bar C_{11}\pm\bar C_{12}\ge m$, $\bar C_{33}\ge m$ ($\Cb\succ0$),
$\hat C_{11}\pm\hat C_{12}\ge m$, $\hat C_{33}\ge m$ ($\Ch\succ0$), $G\ge0$. Columns are scaled as
in Step 11.1 and constraint rows normalized to unit length.

\subsection*{Step 12.2 -- Exact solution by active-set enumeration}
If the unconstrained minimizer $\bth^\ast=\bP^+\bb$ is feasible it is the solution. Otherwise,
since \eqref{eq:qp} is a convex QP, its optimum is the equality-constrained minimizer of its own
active set $\mathcal A$; for every $\mathcal A\subseteq\{1,\dots,n_c\}$ the KKT system
\begin{equation}
  \begin{bmatrix}\bP^T\bP&\bm G_{\mathcal A}^T\\\bm G_{\mathcal A}&\bm0\end{bmatrix}
  \begin{bmatrix}\bth\\\bm\lambda\end{bmatrix}=
  \begin{bmatrix}\bP^T\bb\\\bm h_{\mathcal A}\end{bmatrix}
\end{equation}
is solved (pseudo-inverse), and the feasible candidate with the smallest $J$ is the global optimum
($2^{n_c}-1\le511$ systems). Parameters seen by no equation are reported as ``no information''.
Optionally (\code{requireEnergyPD}) $\lambda_{\min}(\B)\ge m$ is added and solved by \code{fmincon}.

\subsection*{Step 12.3 -- Trend and $N$-dependence of the free values}
Each free value is fitted by $\theta_j(N)\approx c_0+c_1N^{-1}\,(+c_2N^{-2}$ with $\ge5$ values
of $N)$, and its scaling exponent is estimated by
\begin{equation}
  (s_j,b_j)=\arg\min\sum_N\big(\log\theta_j(N)-s\log N-b\big)^2,
\end{equation}
$s_j\approx0$ indicating an $N$-independent (material) constant.
\begin{delivers}
Free $\bth(N)$ per $N$, the global $\bth$, trends and exponents $s_j$ (printed next to the monotone
result for comparison).
\end{delivers}

% ==================================================================
\section{Stage 13 -- Outputs and diagnostic plots}
\label{sec:stage13}
% ==================================================================

\subsection*{Step 13.1 -- Parameter and fit plots}
\implementedin{\code{plotMonotone}, \code{plotFit}}
$\Cb(N)/\C$, $\Ch(N)/\C$ (componentwise), $G(N)$, and $\lambda_{\min}$ of $\Cb$, $\Ch$, $\B$ versus
$N$; per channel the scaled residuals versus $N$ and model--FE scatter plots of the stress and
double-stress components.

\subsection*{Step 13.2 -- Data check: $A_1\dots A_{19}$ per load case}
\implementedin{\code{plotAterms} (one PNG per load case and region)}
Each panel shows raw and smoothed $A_j$ versus $N$, normalized by the macroscopic reference of the
same component pair:
\begin{equation}
  \frac{A_j}{A^M_{\rho(j)}},\qquad
  \rho=(1,2,3,4,\ 1,2,3,4,\ 1,2,2,3,4,\ 14,\dots,19),
\end{equation}
e.g.\ $A_1/\int\varepsilon_{M11}^2=C_{11}A_1/\int\varepsilon_{M11}C_{11}\varepsilon_{M11}$. For
$\bpsi=\epsM$ the ratio is $1$ for $j\le4$ and $j\ge14$ and $0$ for the $\gamma$-terms. The ratio
is used only where the reference is physically relevant,
\begin{equation}
  \frac{c^{\mathrm{ref}}_{\rho(j)}\,|A^M_{\rho(j)}|}{2U_{\mathrm{hom}}}\ge10^{-2}\ (j\le13),\qquad
  \frac{\ell^2C_{11}A^M_j}{2U_{\mathrm{hom}}}\ge10^{-6}\ \text{and}\ A^M_j\ge10^{-3}\!\sum_{i=14}^{19}\!A^M_i\ (j\ge14),
\end{equation}
$c^{\mathrm{ref}}=(C_{11},2|C_{12}|,C_{22},4C_{33})$; otherwise (e.g.\ $\varepsilon_{M11}\equiv0$ under
$H_{22}$, $\varepsilon_{M22}$ a mere equilibrium correction under $G_{1,22}$, $\bkap_M\equiv\bm0$ under
pure $\bH$) the energy fraction $c_jA_j/2U_{\mathrm{hom}}$ is shown, $c_j$ being the weight of
$A_j$ in \eqref{eq:Wcomp} with $\Cb=\Ch=\C$ (curvature terms: $\ell^2C_{11}$).

\noindent Expected behaviour for $N\to\infty$: $\gamma\gamma$-terms decrease to $0$; the
$\psi\gamma$-terms tend to $0$ in magnitude (negative values rise to $0$); the $\psi\psi$-ratios tend
to $1$ from either side; the $\kappa\kappa$-ratios need not be monotone ($\int\kappa^2$ of a boundary
layer grows $\propto N$), only $\ell^2A_j$ decreases.
\begin{delivers}
\code{Params\_vsN\_monotone}, \code{Fit\_monotone}, and
\code{<region>/A\_terms/<case>\_A.png} for every load case.
\end{delivers}

% ==================================================================
\section{Stage 14 -- Consistency properties and checks}
\label{sec:stage14}
% ==================================================================

\begin{proposition}[Vanishing integrals by symmetry]
\label{prop:sym}
Let the region be centred at $\bX_c$ and $\Zset$ symmetric. Then $\bm S_{\mathrm{FE}}=\bm0$ for pure
second-gradient loads ($\bH=\bm0$) and $\bm M_{\mathrm{FE}}=\bm0$ for pure first-gradient loads
($\bG=\bm0$).
\end{proposition}
\begin{proof}
The point reflection $\bx\mapsto-\bx$ maps the microstructure of shift $\bm\zeta$ onto that of
$-\bm\zeta$ (the centres \eqref{eq:centres} form a symmetric set). For $\bH=\bm0$ the data
\eqref{eq:kubc} are even, so $\bu^{(-\bm\zeta)}(\bx)=\bu^{(\bm\zeta)}(-\bx)$ and
$\bsig^{(-\bm\zeta)}(\bx)=-\bsig^{(\bm\zeta)}(-\bx)$: the pair's stress integrals cancel. For
$\bG=\bm0$ the data are odd, $\bsig^{(-\bm\zeta)}(\bx)=\bsig^{(\bm\zeta)}(-\bx)$, and the odd
moment arm cancels the first moments.
\end{proof}
\noindent Hence the stress channel is informative only for load cases with $\bH\ne\bm0$ and the
double-stress channel only for $\bG\ne\bm0$ (criterion \eqref{eq:info}); this is why the library
\code{H3G16} contains the three $H$ cases.

\begin{proposition}[Homogeneous material]
\label{prop:hom}
If $\Dm\equiv\C$, then $\bpsi=\epsM$, $\bgam=\bm0$, $\bkap=\bkap_M$, and $\Ch=\C$, $G=0$ satisfy
\eqref{eq:eqE}--\eqref{eq:eqM} exactly for any $\Cb$ ($\Cb$ is not identifiable).
\end{proposition}
\begin{proof}
All $A_5\dots A_{13}$, $I_4\dots I_6$, $I_{19}\dots I_{24}$ vanish; \eqref{eq:eqE} reduces to
$\int\epsM^T\C\epsM=2U_{\mathrm{het}}$, \eqref{eq:eqS} to $\C\int\ve{\epsM}=\bm S_{\mathrm{FE}}$ and
\eqref{eq:eqM} to $\bm D\int\bm m(\C\ve{\epsM},\bx)=\bm D\bm M_{\mathrm{FE}}$ with $G=0$.
\end{proof}

\paragraph{Numerical checks performed by the code.}
\begin{itemize}[leftmargin=1.2em]
\item PBC: recovered average strain equals the prescribed one ($<10^{-8}$).
\item Quadrature: $\int\epsM^T\C\epsM$ (grid) vs.\ $2U_{\mathrm{hom}}$ (FE), $\sim10^{-3}$.
\item Local gradient exact for affine fields (Lemma~\ref{lem:mls}); region weights reproduce the
trapezoidal rule and exact box areas.
\item Monotone solution: $\lambda_{\min}(\B)\ge m$ at every $N$, coefficients $\ge0$.
\end{itemize}

% ==================================================================
\appendix
\section{Default settings}
% ==================================================================
\begin{center}\small
\begin{longtable}{@{}lll@{}}
\toprule
symbol / field & default & meaning\\
\midrule
$L_x=L_y$ & $1$ & specimen size\\
$\mathcal N$ (\code{Nlist}) & $1,\dots,8$ & unit cells per side\\
$r$ (\code{Rfrac}) & $6/19$ & inclusion radius / cell size\\
$h/\ell$, $h_{\C}/L$ & $0.05$, $0.02$ & element size (heterogeneous; Stage 1 and PBC)\\
$E_m,\nu_m$; $E_i,\nu_i$ & $70000, 0.33$; $3500, 0.33$ & phases (plane strain)\\
$\varepsilon_0$, $g$ & $0.01$, $2$ & load amplitude, gradient scale\\
library & \code{H3G16} & 19 load cases\\
$n$ (\code{nPeriod}) & $8$ & shifts per direction, $n^2$ solves per $N$\\
$n_g$, $n_{\mathrm{cell}}$ & $121$, $40$ & grid points per side (minimum, per cell)\\
\code{kappaMethod}, $w_\sigma$ & \code{'local'}, $0.1$ & $\nabla\epsM$ exact $-$ MLS $\nabla\bgam$, $\sigma=w_\sigma\ell$\\
$\beta$ (\code{interiorBand}) & $0.5$ & interior band in cells\\
regions & full, interior & identified separately\\
$N_{\min}$ & $1$ & smallest $N$ used\\
$\mathcal P_d$, $\mathcal P_f$ & $\{\tfrac12,1,\tfrac32,2\}$, $\{0,1,2\}$ & smoothing exponents\\
mode & \code{'monotone'} & Stage 11 (per-$N$ free fit for comparison)\\
\code{Gmode} & \code{'single'} & $G_{11}=G_{22}=G_{33}$\\
$\bm e$ & $(0,1,2)$ & exponents of $\theta_j(N)$\\
$w_E,w_S,w_D$ & $1,1,1$ & channel weights\\
\code{doubleStressModel} & \code{'local+mu'} & \eqref{eq:eqM}\\
$\tau$ & $10^{-2}$ & informativeness threshold\\
$\delta$ (\code{pdMargin}) & $10^{-3}$ & positive-definiteness margin\\
\code{energyPD} & true & $\B\succ0$ at every $N$\\
\bottomrule
\end{longtable}
\end{center}

% ==================================================================
\section{Algorithm}
% ==================================================================
\begin{algorithm}[h]
\caption{Identification of $\Cb(N)$, $\Ch(N)$, $G(N)$}
\begin{algorithmic}[1]
\Require load cases $\Kset$, cell numbers $\Nset$, material, settings
\State $\C\gets$ periodic homogenization of one unit cell \Comment{Stage 2}
\State solve the plain domain with $\C$ under KUBC $\to\epsM^{(k)}$, $2U^{(k)}_{\mathrm{hom}}$ \Comment{Stage 5}
\For{$N\in\Nset$}
  \For{$\bm\zeta\in\Zset$} \Comment{Stage 6}
    \State mesh, assemble $\bK$, solve all $k$; accumulate $\eps$ on the grid, $\bu^T\bK_r\bu$, $\bm Q_s\bu$, $\bm Q_m\bu$
  \EndFor
  \State $\bpsi^{(k)}$, $2U_{\mathrm{het}}$, $\bm S_{\mathrm{FE}}$, $\bm M_{\mathrm{FE}}$ (both regions)
  \State $\bgam=\epsM-\bpsi$, $\bkap=\nabla\epsM-\nabla_\sigma\bgam$, $\bkap_M=\nabla\epsM$ \Comment{Stage 7}
  \State $A_{1..19}$, $I_{1..24}$, $A^M$ for both regions \Comment{Stage 8}
\EndFor
\For{region $\in\{\text{full},\text{interior}\}$}
  \State smooth every density over $N$ (Table~\ref{tab:smooth}) \Comment{Stage 9}
  \State assemble \eqref{eq:eqE}--\eqref{eq:eqM} for all informative groups \Comment{Stage 10}
  \State solve \eqref{eq:monoopt} $\to$ $c_{jk}$, $\bth(N)$ \Comment{Stage 11}
  \State free per-$N$ and global solutions, exponents $s_j$ \Comment{Stage 12}
  \State plots and one PNG per load case \Comment{Stage 13}
\EndFor
\Ensure $\bth(N)=[\bar C_{11},\bar C_{12},\bar C_{33},\hat C_{11},\hat C_{12},\hat C_{33},G](N)$ per region
\end{algorithmic}
\end{algorithm}

\end{document}
